% ---------------------------------------------------------------------
%  Abstract
% ---------------------------------------------------------------------
\chapter*{Abstract}
\addcontentsline{toc}{chapter}{Abstract}
\markboth{Abstract}{}

A system can tip into a new state not only because a parameter crosses a critical value, but also
because the parameter changes faster than the system can follow. This second route, \emph{rate-induced
tipping}, is well developed in theory and in models, yet laboratory demonstrations remain few. In this
project we ask whether a simple electronic circuit can show it.

We study an operational-amplifier Wien-bridge oscillator whose gain $G$ is swept back and forth
across the oscillation threshold $\Gc$ by an Arduino-driven digital potentiometer. Above $\Gc$ the quiet
state loses stability through a Hopf bifurcation and an oscillation grows out of noise; below $\Gc$ it
decays. The work proceeds in four stages. Circuit simulations of a smooth gain ramp first show that the
onset and the decay of the oscillation both lag behind the threshold, and that the lag grows as the
square root of the sweep rate $\lambda$. Bench experiments with a staircase sweep of the feedback
resistor confirm the same $\sqrt{\lambda}$ law on hardware and reveal rate-dependent hysteresis. A
simulated copy of that experiment extends the sweep to rates the bench cannot reach. Finally, sweeping
the input resistor along a fixed gain path, we find that the long-time state of the circuit depends on
the rate alone: at slow rates the oscillation re-ignites in every cycle, while above a sharp critical
rate it switches off once and never returns.

A single growth-rate argument accounts for all of these observations. The oscillation needs a fixed
amount of exponential growth to climb from the noise floor to full size, and a faster sweep leaves less
time above threshold in which to collect it. The same argument yields the $\sqrt{\lambda}$ law, the
critical rates, and their dependence on noise. Together, the results establish the Wien-bridge
oscillator as a controllable analogue system that displays clear signatures of rate-induced tipping.
